\documentclass[11pt]{article}
\usepackage{mystyle}

\title{Rosen --- \S 5.1: Mathematical Induction\\ \large Selected Exercises}
\author{Alston}
\date{Semester 2, 2026}

\begin{document}
\maketitle

% ======================================================================
\begin{exercise}{25}
    Prove that if $h > -1$, then $1 + nh \leq (1+h)^n$ for all
    nonnegative integers $n$. This is called \emph{Bernoulli's
    inequality}.
\end{exercise}

% ======================================================================
\begin{exercise}{27}
    Prove that for every positive integer $n$,
    \[
    1 + \frac{1}{\sqrt{2}} + \frac{1}{\sqrt{3}} + \cdots +
    \frac{1}{\sqrt{n}} > 2\big(\sqrt{n+1} - 1\big).
    \]
\end{exercise}

% ======================================================================
\begin{exercise}{29}
    \textit{In this and the next exercise, $H_n$ denotes the $n$th
    harmonic number, $H_n = 1 + \frac12 + \frac13 + \cdots +
    \frac1n$.}

    Prove that $H_{2^n} \leq 1 + n$ whenever $n$ is a nonnegative
    integer.
\end{exercise}

% ======================================================================
\begin{exercise}{30}
    Prove that
    \[
    H_1 + H_2 + \cdots + H_n = (n+1)H_n - n.
    \]
\end{exercise}

% ======================================================================
\begin{exercise}{35}
    Prove that $n^2 - 1$ is divisible by $8$ whenever $n$ is an odd
    positive integer.
\end{exercise}

% ======================================================================
\begin{exercise}{37}
    Prove that if $n$ is a positive integer, then $133$ divides
    $11^{n+1} + 12^{2n-1}$.
\end{exercise}

% ======================================================================
\begin{exercise}{40}
    Prove that if $A_1, A_2, \dots, A_n$ and $B$ are sets, then
    \[
    (A_1 \cap A_2 \cap \cdots \cap A_n) \cup B
    = (A_1 \cup B) \cap (A_2 \cup B) \cap \cdots \cap (A_n \cup B).
    \]
\end{exercise}

% ======================================================================
\begin{exercise}{46}
    Prove that a set with $n$ elements has $n(n-1)(n-2)/6$ subsets
    containing exactly three elements whenever $n$ is an integer
    greater than or equal to $3$.
\end{exercise}

% ======================================================================
\begin{exercise}{52}
    Suppose that $m$ and $n$ are positive integers with $m > n$ and
    $f$ is a function from $\{1, 2, \dots, m\}$ to $\{1, 2, \dots, n\}$.
    Use mathematical induction on the variable $n$ to show that $f$ is
    not one-to-one.
\end{exercise}

% ======================================================================
\begin{exercise}{53}
    Use mathematical induction to show that $n$ people can divide a
    cake (where each person gets one or more separate pieces of the
    cake) so that the cake is divided fairly, that is, in the sense
    that each person thinks he or she got at least $(1/n)$th of the
    cake.
    [\emph{Hint}: For the inductive step, take a fair division of the
    cake among the first $k$ people, have each person divide their
    share into what this person thinks are $k+1$ equal portions, and
    then have the $(k+1)$st person select a portion from each of the
    $k$ people. When showing this produces a fair division for $k+1$
    people, suppose that person $k+1$ thinks that person $i$ got
    $p_i$ of the cake where $\sum_{i=1}^{k} p_i = 1$.]
\end{exercise}

% ======================================================================
\begin{exercise}{55}
    A knight on a chessboard can move one space horizontally (in
    either direction) and two spaces vertically (in either direction)
    or two spaces horizontally (in either direction) and one space
    vertically (in either direction). Suppose that we have an infinite
    chessboard, made up of all squares $(m, n)$ where $m$ and $n$ are
    nonnegative integers that denote the row number and the column
    number of the square, respectively. Use mathematical induction to
    show that a knight starting at $(0,0)$ can visit every square
    using a finite sequence of moves.
    [\emph{Hint}: Use induction on the variable $s = m + n$.]
\end{exercise}

% ======================================================================
\begin{exercise}{61}
    Show that
    \[
    \big[(p_1 \to p_2) \wedge (p_2 \to p_3) \wedge \cdots \wedge
    (p_{n-1} \to p_n)\big]
    \to \big[(p_1 \wedge p_2 \wedge \cdots \wedge p_{n-1}) \to p_n\big]
    \]
    is a tautology whenever $p_1, p_2, \dots, p_n$ are propositions,
    where $n \geq 2$.
\end{exercise}

% ======================================================================
\begin{exercise}{62}
    Show that $n$ lines separate the plane into $(n^2+n+2)/2$ regions
    if no two of these lines are parallel and no three pass through a
    common point.
\end{exercise}

% ======================================================================
\begin{exercise}{63}
    Let $a_1, a_2, \dots, a_n$ be positive real numbers. The
    \emph{arithmetic mean} of these numbers is defined by
    $A = (a_1 + a_2 + \cdots + a_n)/n$, and the \emph{geometric mean}
    of these numbers is defined by $G = (a_1 a_2 \cdots a_n)^{1/n}$.
    Use mathematical induction to prove that $A \geq G$.
\end{exercise}

% ======================================================================
\begin{exercise}{65}
    Show that if $n$ is a positive integer, then
    \[
    \sum_{\{a_1, \dots, a_k\} \subseteq \{1, 2, \dots, n\}}
    \frac{1}{a_1 a_2 \cdots a_k} = n.
    \]
    (Here the sum is over all nonempty subsets of the set of the $n$
    smallest positive integers.)
\end{exercise}

% ======================================================================
\begin{exercise}{68}
    A guest at a party is a \emph{celebrity} if this person is known
    by every other guest, but knows none of them. There is at most
    one celebrity at a party, for if there were two, they would know
    each other. A particular party may have no celebrity. Your
    assignment is to find the celebrity, if one exists, at a party, by
    asking only one type of question---asking a guest whether they
    know a second guest. Everyone must answer your questions
    truthfully. That is, if Alice and Bob are two people at the
    party, you can ask Alice whether she knows Bob; she must answer
    correctly.

    Use mathematical induction to show that if there are $n$ people at
    the party, then you can find the celebrity, if there is one, with
    $3(n-1)$ questions.
    [\emph{Hint}: First ask a question to eliminate one person as a
    celebrity. Then use the inductive hypothesis to identify a
    potential celebrity. Finally, ask two more questions to determine
    whether that person is actually a celebrity.]
\end{exercise}

% ======================================================================
\begin{exercise}{71}
    \textit{Suppose there are $n$ people in a group, each aware of a
    scandal no one else in the group knows about. These people
    communicate by telephone; when two people in the group talk, they
    share information about all scandals each knows about. The
    gossip problem asks for $G(n)$, the minimum number of telephone
    calls that are needed for all $n$ people to learn about all the
    scandals.}

    Prove that $G(n) = 2n - 4$ for $n \geq 4$.
\end{exercise}

% ======================================================================
\begin{exercise}{73}
    Show that if $I_1, I_2, \dots, I_n$ is a collection of open
    intervals on the real number line, $n \geq 2$, and every pair of
    these intervals has a nonempty intersection, that is,
    $I_i \cap I_j \neq \emptyset$ whenever $1 \leq i \leq n$ and
    $1 \leq j \leq n$, then the intersection of all these sets is
    nonempty, that is, $I_1 \cap I_2 \cap \cdots \cap I_n \neq
    \emptyset$.
\end{exercise}

\end{document}